\scanpage{10}
ou encore
\[
\exists\ q,r\in\hat{\mathbb Z}^{\unsure{*}}\quad\text{tels que}\quad\boxed{gl_1^q\sigma_\infty(g)=l_0^r}
\]
\margin{Il vaut mieux écrire la relation sous la forme
\[
\sigma_\infty(g^{-1})=l_0^\beta gl_1^\gamma,
\]
on trouve, en \unsure{passant} au groupe \unsure{abélien}, que $\beta=\gamma$ (i.e. $q=-r$) donc
\[
\sigma_\infty(g^{-1})=l_0^\beta gl_1^\beta
\]}
(Condition qui ne dépend que de $g$, pas de $p$;

Si $g$ est remplacé par $g'=l_0^\alpha gl_1^{-\alpha}$, on aura $\sigma_\infty(g')=l_1^\alpha\sigma_\infty(g)l_0^{-\alpha}$, et la relation
\[
g'l_1^{q'}\sigma_\infty(g')=l_0^{r'}
\]
s'écrit
\[
l_0^\alpha gl_1^{-\alpha}l_1^{q'}l_1^\alpha\sigma_\infty(g)l_0^{-\alpha}=l_0^{r'}
\]
i.e. en simplifiant
\[
gl_1^{q'}\sigma_\infty(g)=l_0^{r'}
\]
donc $q=q'$, $r=r'$ i.e. $(q,r)$ ne dépendent que de l'autom ext. défini par $u$, et pas du \deleted{\unsure{choix de u}} de son représentant dans le normalisateur de $L_0$.

Dans le cas de $\tau$, où $g=l_1$, la relation devient
\[
l_1l_1^ql_0=l_0^r\quad\text{donc}\quad q=-1,\ r=1.
\]
On sait que, si on travaille sur les groupes discrets, en plus de l'autom ext. identique, il n'y a qu'une seule autre solution de nos équations (en exigeant \inserted{donc $p=\pm1$ et} que $l_0^p$ et $\operatorname{int}(g)l_1^p$ engendrent $\pi$, en plus des relations fondamentales
\mspage{9}
